\documentclass[a4paper]{article}

% Use Package
\usepackage[english]{babel}
\usepackage[utf8]{inputenc}
\usepackage{amsmath}
\usepackage{graphicx}
\usepackage[colorinlistoftodos]{todonotes}
\usepackage{fontspec}
\usepackage{geometry}
\usepackage{layout}
\usepackage{titling}
\usepackage{listings}  

% Define Commands
\newcommand{\subtitle}[1]{%
  \posttitle{%
    \par\end{center}
    \begin{center}\large#1\end{center}
    \vskip0.5em}%
}

% Initial Set Up
\setmainfont{Arial}
\lstdefinestyle{codeStyle}{
  belowcaptionskip=1\baselineskip,
  breaklines=true,
  frame=L,
  xleftmargin=\parindent,
  language=Java,
  showstringspaces=false,
  basicstyle=\footnotesize\ttfamily,
  keywordstyle=\bfseries\color{green!40!black},
  commentstyle=\itshape\color{purple!40!black},
  identifierstyle=\color{blue},
  stringstyle=\color{orange},
}
\lstset{ 
  backgroundcolor=\color{white}, 
  breakatwhitespace=false,         
  breaklines=true,              
  captionpos=b,                    
  commentstyle=\color{green},    
  deletekeywords={...},            
  escapeinside={\%*}{*)},          
  extendedchars=true,              
  frame=single,	                 
  keepspaces=true,                
  keywordstyle=\color{blue},      
  language=Java,                  
  otherkeywords={*,...},           
  numbers=left,                    
  numbersep=5pt,                   
  numberstyle=\color{gray},
  rulecolor=\color{white},         
  showspaces=false,             
  showstringspaces=false,          
  showtabs=false,                  
  stepnumber=1,                    
  stringstyle=\color{mauve},     
  tabsize=1,	                   
  title=\lstname,                   
  escapechar=@,
  style=codeStyle
}

\title{Operators}
\subtitle{Java Fundamentals}
\date{} 
\author{Libre Education}
\begin{document}
\maketitle


\begin{figure}[b]
\includegraphics{BY-SA}
\centering
\end{figure}

\newpage

\section*{Operators}

An operator is a construct that performs a function to one or more values.

\section*{Types of Operators}
Operators are grouped by what kind of operation they perform.

\subsection*{Arithmetic Operators}
Arithmetic operators perform typical mathematical functions to two values.
\begin{enumerate}

\item = (Simple Assignment) - Assigns the value on the right to the value on the left

\item += (Add And Assign) - Adds the two values together and then assigns the value on the left to the result

\item -= (Subtract And Assign) - Subtracts the right value from the left value and then assigns the value on the left to the result

\item *= (Multiply And Assign) - Multiplies the values together and then assigns the value on the left to the result

\item /= (Divide And Assign) - Divided the value on the left by the value on the right and then assigns the value on the left to the result

\end{enumerate}

\subsection*{Unary Operators}
Unary operators require and perform a function to one value.

\begin{enumerate}

\item + (Unary Plus)

\item - (Unary Minus)

\item + + (Increment) - Increases the value by 1

\item - - (Decrement) - Decreases the value by 1

\item ! (Logical Compliment) - Inverts the value of a boolean

\end{enumerate}

\subsection*{Equality And Relational Operators}
Equality and Relational operators check for a condition and return a boolean but do not manipulate values
\begin{enumerate}

\item == (Equal To) - Returns true if the two values are equal to each other

\item != (Not Equal To) - Returns true if the two values are not equal to each other

\item > (Greater Than) - Returns true if the value on the left is greater than the value on the right

\item < (Less Than) - Returns true if the value on the left is less than the value on the right

\item >= (Greater Than or Equal To) - Returns true if the value on the left is greater than or equal to the value on the right

\item <= (Less Than or Equal To) - Returns true of the value on the left is less than or equal to the value on the right

\item instanceof (Instance of) - Compares an object and a class. Returns true if the object provided on the left is an instance of the class provided on the right

\end{enumerate}

\subsection*{Conditional Operators}
Conditional operators check two statements or expressions and return a boolean based on the result
\begin{enumerate}

\item \&\& (Conditional AND) - Returns true if both statements are true

\item || (Conditional OR) - Returns true if at least one of the statements are true

\end{enumerate}

\subsection*{Bitwise and Bitshift Operators}
Bitwise and Bitshift operators manipulate values at the byte level. They take place between each parallel pair of bits in the value
\begin{enumerate}

\item $\sim$ (Unary Bitwise Complement) - Inverts each bit in the value

\item \& (Bitwise AND) - If both corresponding bits are 1, the result is 1

\item | (Bitwise Inclusive OR) - If either of the corresponding bits are 1, the result is 1

\item \string^ (Bitwise Exclusive OR) - If both of the corresponding bits are 1, the result is 1

\end{enumerate}

\subsection*{Ternary Operator}
? : - The ternary operator provides a shorter way to write and if...then...else statement and takes three operands: an expression and two values. If the expression evaluates to true the first value is returned and if the expression evaluates to false the second value is returned.

\subsubsection*{Syntax}
\begin{lstlisting}
result = expression ? value1 : value2;
\end{lstlisting}

\subsubsection*{Example}
\begin{lstlisting}
boolean expression = true;
String result = expression ? "The expression is true" : "The expression is false";
System.out.println(result);
// Output: The expression is true

\end{lstlisting}


\newpage

\section*{Other Resources}
\begin{enumerate}

\item Wikipedia (en.wikipedia.org/wiki/Operator\_(computer\_programming))

\item The Java Tutorials (docs.oracle.com/javase/tutorial/java/nutsandbolts/operators.html)

\item TutorialsPoint (www.tutorialspoint.com/java/java\_basic\_operators.htm) \\ 

Bitwise and Bitshift Operators

\item Rose India (www.roseindia.net/java/master-java/bitwise-bitshift-operators.shtml)

\end{enumerate}


\end{document}